The angles $\angle AXD$ and $\angle BXC$ are equal as vertical angles. Also $\dfrac{XA}{XB}=\dfrac{XD}{XC}=\dfrac{1}{2}$. The cosine rule, or side-angle-side similarity with the two rays interchanged, therefore gives
\[
AD=\dfrac{1}{2}BC.
\]
Using the two cyclic quadrilaterals $AXDY$ and $B X C Y$, and the collinearities $A,X,C$ and $B,X,D$, inscribed-angle chasing gives
\[
\angle YAD=\angle YCB,\qquad \angle YDA=\angle YBC.
\]
Consequently triangles $YAD$ and $YCB$ are similar. Their corresponding sides satisfy
\[
\dfrac{YD}{YB}=\dfrac{AD}{CB}=\dfrac{1}{2}.
\]
Thus $\boxed{YD=6}$.
\begin{asy}
texpreamble("\usepackage{amsmath,amssymb}");
size(240);
defaultpen(fontsize(9pt));
real c=-571.0/760,s=sqrt(1-c*c);pair X=(0,0),A=(4,0),B=8(c,s),C=(-14,0),D=-7(c,s),O1=(2,(-7-4c)/(2s)),O2=(-7,(8+14c)/(2s)),Y=-22*(5+6c,6s)/(61+60c);draw(circle(O1,length(O1)),gray);draw(circle(O2,length(O2)),gray);draw(A--B--C--D--cycle);draw(A--C,dashed);draw(B--D,dashed);draw(A--Y--D,blue);draw(C--Y--B,red);label("$A$",A,(1,0));label("$B$",B,(0,1));label("$C$",C,(-1,0));label("$D$",D,SE);label("$X$",X,(0,1));label("$Y$",Y,(0,-1));
\end{asy}
